% !TeX program = xelatex
\documentclass[aspectratio=169,10pt]{beamer}

\usetheme{default}
\usefonttheme[onlymath]{serif}

\usepackage{amsmath,amssymb}
\usepackage{booktabs}
\usepackage{tabularx}
\usepackage{tikz}
\usepackage{iftex}
\usetikzlibrary{positioning,arrows.meta,calc}

\ifPDFTeX
  \usepackage[utf8]{inputenc}
  \usepackage[T1]{fontenc}
\else
  \usepackage{fontspec}
  \usepackage{xeCJK}
  \IfFontExistsTF{Arial}{\setsansfont{Arial}}{\setsansfont{TeX Gyre Heros}}
  \IfFontExistsTF{Times New Roman}{\setmainfont{Times New Roman}}{\setmainfont{TeX Gyre Termes}}
  \IfFontExistsTF{Heiti SC}{%
    \setCJKmainfont{Heiti SC}
    \setCJKsansfont{Heiti SC}
  }{%
    \IfFontExistsTF{FandolHei-Regular.otf}{%
      \setCJKmainfont[BoldFont=FandolHei-Regular.otf]{FandolSong-Regular.otf}
      \setCJKsansfont[BoldFont=FandolHei-Regular.otf]{FandolHei-Regular.otf}
    }{%
      \IfFontExistsTF{Songti SC}{%
        \setCJKmainfont{Songti SC}
        \setCJKsansfont{Songti SC}
      }{%
        \setCJKmainfont{STSong}
        \setCJKsansfont{STSong}
      }
    }
  }
\fi
\renewcommand{\familydefault}{\sfdefault}

\definecolor{BrandOrange}{RGB}{243,93,0}
\definecolor{BrandOrangeDark}{RGB}{184,67,0}
\definecolor{BrandOrangeLight}{RGB}{255,239,232}
\definecolor{BrandInk}{RGB}{34,34,34}
\definecolor{BrandGray}{RGB}{91,91,91}
\definecolor{BrandBlue}{RGB}{28,76,132}
\definecolor{BrandBlueLight}{RGB}{226,239,255}
\definecolor{BrandGreen}{RGB}{34,124,86}
\definecolor{BrandGreenLight}{RGB}{229,246,237}
\definecolor{BrandPurple}{RGB}{96,71,143}
\definecolor{BrandPurpleLight}{RGB}{239,233,249}
\definecolor{BrandRed}{RGB}{178,55,55}
\definecolor{BrandRedLight}{RGB}{255,235,235}

\setbeamertemplate{navigation symbols}{}
\setbeamercolor{normal text}{fg=BrandInk,bg=white}
\setbeamercolor{structure}{fg=BrandOrange}
\setbeamercolor{block title}{fg=white,bg=BrandOrange}
\setbeamercolor{block body}{fg=BrandInk,bg=BrandOrangeLight}
\setbeamertemplate{itemize item}{\color{BrandOrange}$\blacktriangleright$}
\setlength{\parskip}{0.10em}
\setlength{\tabcolsep}{4pt}
\renewcommand{\arraystretch}{1.10}

\newcommand{\tightitems}{\setlength{\itemsep}{0.20em}}
\newcommand{\key}[1]{\textcolor{BrandOrangeDark}{\textbf{#1}}}
\newcommand{\muted}[1]{\textcolor{BrandGray}{#1}}
\newcommand{\contentheader}[1]{%
  \vspace*{0.06cm}%
  {\noindent\color{BrandOrange}\rule{0.42em}{0.42em}\hspace{0.50em}\bfseries\Large #1\par}%
  \vspace{0.10em}{\color{BrandOrange}\rule{\textwidth}{1.25pt}\par}\vspace{0.25em}%
}
\newcommand{\tagbox}[4]{%
  \node[draw=#1,fill=#2,rounded corners=2pt,text width=#3,align=center,inner sep=5pt,font=\scriptsize] {#4};
}

\title{OpenPEAK：金融内网可治理 Agent 底座}
\subtitle{从“能回答”到“可控、可查、可复用”的智能工作流}
\author{Yida Ding}
\date{June 2026}

\begin{document}
\small

\begin{frame}[plain]
\begin{tikzpicture}[remember picture,overlay]
  \fill[BrandOrange] (current page.south west) rectangle (current page.north east);
  \fill[BrandOrangeDark] (current page.south west) rectangle ([xshift=0.18\paperwidth]current page.north west);
  \foreach \x in {0.24,0.34,0.46,0.59,0.73,0.86} {
    \draw[white!20,line width=0.5pt] ([xshift=\x\paperwidth]current page.south west) -- ([xshift=\x\paperwidth,yshift=0.38\paperheight]current page.south west);
  }
  \node[align=left,text=white,font=\bfseries\fontsize{25}{30}\selectfont,text width=10.5cm] at ([xshift=1.8cm,yshift=0.82cm]current page.center) {OpenPEAK};
  \node[align=left,text=white,font=\bfseries\fontsize{18}{24}\selectfont,text width=10.5cm] at ([xshift=1.8cm,yshift=-0.05cm]current page.center) {金融内网可治理 Agent 底座};
  \node[align=left,text=white!88,font=\large,text width=10.5cm] at ([xshift=1.8cm,yshift=-0.92cm]current page.center) {可控、可查、可复用：把大模型变成交易相关场景可接受的工作流系统};
  \node[align=left,text=white!82,font=\normalsize,text width=10.5cm] at ([xshift=1.8cm,yshift=-1.82cm]current page.center) {Yida Ding \quad | \quad June 2026};
  \node[draw=white,fill=white,rounded corners=2pt,text=BrandOrangeDark,font=\bfseries\small,inner sep=7pt] at ([xshift=-1.25cm,yshift=-2.82cm]current page.center) {Direct 轻量问答 \quad + \quad OGF Expert 专家流程 \quad + \quad LLM Wiki 本地沉淀};
\end{tikzpicture}
\end{frame}

\begin{frame}[plain]
\contentheader{为什么领导需要关心：裸 LLM 进不了金融交易流程}
\begin{columns}[T,totalwidth=\textwidth]
\begin{column}{0.47\textwidth}
\begin{tikzpicture}[x=1cm,y=1cm,every node/.style={font=\scriptsize}]
  \node[draw=BrandRed,fill=BrandRedLight,rounded corners=2pt,text width=5.2cm,align=left,inner sep=6pt] (a) at (2.8,4.05) {\key{不能外发数据}\\交易、风控、研究材料无法依赖公网搜索};
  \node[draw=BrandOrange,fill=BrandOrangeLight,rounded corners=2pt,text width=5.2cm,align=left,inner sep=6pt] (b) at (2.8,2.75) {\key{答案难审计}\\用了什么上下文、为什么这么判断，不易复盘};
  \node[draw=BrandPurple,fill=BrandPurpleLight,rounded corners=2pt,text width=5.2cm,align=left,inner sep=6pt] (c) at (2.8,1.45) {\key{知识不沉淀}\\一次排查结束后，经验没有变成组织资产};
  \draw[-{Stealth[length=2mm]},BrandGray,thick] (a) -- (b);
  \draw[-{Stealth[length=2mm]},BrandGray,thick] (b) -- (c);
\end{tikzpicture}
\end{column}
\begin{column}{0.50\textwidth}
\begin{block}{管理层真正要买的是治理能力}
\begin{itemize}\tightitems
  \item 能否在内网运行，减少数据外泄风险？
  \item 出错后能否追溯到上下文、日志和证据？
  \item 简单问题能否低成本处理，复杂问题再走专家流程？
  \item 能否把复盘结果留在团队知识库里？
\end{itemize}
\end{block}
\vspace{0.35em}
\begin{center}
{\large\bfseries 普通 AI 助手回答问题；OpenPEAK 管理一次可复盘的工作流。}
\end{center}
\end{column}
\end{columns}
\end{frame}

\begin{frame}[plain]
\contentheader{OpenPEAK 解法：双模式 + OGF 约束 + 本地证据链}
\begin{center}
\begin{tikzpicture}[x=1cm,y=1cm,every node/.style={font=\scriptsize}]
  \node[draw=BrandOrange,fill=BrandOrangeLight,rounded corners=2pt,text width=2.35cm,align=center,inner sep=6pt,font=\scriptsize\bfseries] (user) at (0.95,3.45) {用户入口\\Web / TUI / CLI};
  \node[draw=BrandBlue,fill=BrandBlueLight,rounded corners=2pt,text width=2.35cm,align=center,inner sep=6pt,font=\scriptsize\bfseries] (router) at (3.55,3.45) {模式路由\\显式选择};
  \node[draw=BrandGreen,fill=BrandGreenLight,rounded corners=2pt,text width=2.55cm,align=center,inner sep=6pt,font=\scriptsize\bfseries] (direct) at (6.35,4.45) {Direct\\简单问答\\低 token};
  \node[draw=BrandPurple,fill=BrandPurpleLight,rounded corners=2pt,text width=2.55cm,align=center,inner sep=6pt,font=\scriptsize\bfseries] (ogf) at (6.35,2.45) {OGF Expert\\FICC / 量化\\证据流程};
  \node[draw=BrandGray,fill=white,rounded corners=2pt,text width=2.55cm,align=center,inner sep=6pt,font=\scriptsize\bfseries] (llm) at (9.45,3.45) {内网模型网关\\Local Endpoint};
  \node[draw=BrandOrange,fill=white,rounded corners=2pt,text width=2.05cm,align=center,inner sep=5pt] (logs) at (12.05,4.45) {JSONL Logs\\可追溯};
  \node[draw=BrandOrange,fill=white,rounded corners=2pt,text width=2.05cm,align=center,inner sep=5pt] (wiki) at (12.05,2.45) {LLM Wiki\\知识沉淀};

  \draw[-{Stealth[length=2mm]},BrandOrange,thick] (user.east) -- (router.west);
  \draw[-{Stealth[length=2mm]},BrandGreen,thick] (router.east) -- (direct.west);
  \draw[-{Stealth[length=2mm]},BrandPurple,thick] (router.east) -- (ogf.west);
  \draw[-{Stealth[length=2mm]},BrandGray,thick] (direct.east) -- (llm.west);
  \draw[-{Stealth[length=2mm]},BrandGray,thick] (ogf.east) -- (llm.west);
  \draw[-{Stealth[length=2mm]},BrandOrange,thick] (llm.east) -- (logs.west);
  \draw[-{Stealth[length=2mm]},BrandOrange,thick] (ogf.east) -- ++(1.18,0) |- (wiki.west);

  \node[draw=BrandBlue,fill=white,rounded corners=2pt,text width=11.4cm,align=center,inner sep=5pt,font=\scriptsize\bfseries] at (6.6,0.72) {核心原则：简单问题走 Direct；高风险、专业、需要证据的问题走 OGF Expert。用户必须显式知道当前模式。};
\end{tikzpicture}
\end{center}
\end{frame}

\begin{frame}[plain]
\contentheader{当前可演示能力：已经不是 PPT 原型}
\begin{columns}[T,totalwidth=\textwidth]
\begin{column}{0.55\textwidth}
\begin{tabularx}{\linewidth}{>{\bfseries}p{2.2cm}X}
\toprule
模块 & 已具备能力 \\
\midrule
Web / TUI & 浏览器和终端双入口；默认 OGF，可切 Direct \\
OGF Runtime & ontology / profile / evidence flow，适配 FICC、量化、排查 \\
Doctor & 模型、OGF、Qlib、日志、workspace 一键体检 \\
Trace Logs & 日期分区 JSONL，支持审计、debug 和 Agent 自查 \\
LLM Wiki & 本地 markdown 知识库，沉淀复盘与专家经验 \\
Intranet & 核心不依赖 Node.js；Web 资源嵌入 Go binary \\
\bottomrule
\end{tabularx}
\end{column}
\begin{column}{0.42\textwidth}
\begin{tikzpicture}[x=1cm,y=0.78cm,every node/.style={font=\scriptsize}]
  \draw[->,BrandGray] (0,0) -- (5.35,0);
  \draw[->,BrandGray] (0,0) -- (0,4.8);
  \node[BrandGray,rotate=90] at (-0.38,2.45) {组织价值};
  \fill[BrandOrange] (0.65,0) rectangle (1.05,3.0);
  \fill[BrandBlue] (1.65,0) rectangle (2.05,3.5);
  \fill[BrandPurple] (2.65,0) rectangle (3.05,4.1);
  \fill[BrandGreen] (3.65,0) rectangle (4.05,3.7);
  \fill[BrandOrangeDark] (4.65,0) rectangle (5.05,4.3);
  \node[below,align=center] at (0.85,0) {UI};
  \node[below,align=center] at (1.85,0) {诊断};
  \node[below,align=center] at (2.85,0) {OGF};
  \node[below,align=center] at (3.85,0) {Wiki};
  \node[below,align=center] at (4.85,0) {内网};
  \node[above] at (2.9,4.55) {\bfseries 可进入试点};
\end{tikzpicture}
\vspace{0.3em}
\begin{block}{现场 Demo 顺序}
Readiness 面板 \(\rightarrow\) Direct 简问 \(\rightarrow\) OGF 专家问答 \(\rightarrow\) Trace / Wiki 追溯。
\end{block}
\end{column}
\end{columns}
\end{frame}

\begin{frame}[plain]
\contentheader{建议明天争取的决策：批准 2--4 周受控试点}
\begin{columns}[T,totalwidth=\textwidth]
\begin{column}{0.46\textwidth}
\begin{block}{试点场景}
\begin{itemize}\tightitems
  \item FICC / 量化知识问答与方法复核
  \item 交易、风控、工程问题的证据化排查
  \item 本地 LLM Wiki：把复盘变成团队知识资产
\end{itemize}
\end{block}
\vspace{0.25em}
\begin{block}{需要资源}
\begin{itemize}\tightitems
  \item 一个业务试点团队和 10--20 个黄金问题
  \item 内网模型网关或模型额度
  \item 一批脱敏材料、日志样例、知识库样例
  \item 安全 / 运维 / 业务验收联系人
\end{itemize}
\end{block}
\end{column}
\begin{column}{0.50\textwidth}
\begin{tikzpicture}[x=1cm,y=1cm,every node/.style={font=\scriptsize}]
  \draw[BrandGray,line width=0.8pt] (0.35,2.2) -- (6.2,2.2);
  \foreach \x/\t in {0.35/第1周,2.30/第2周,4.25/第3周,6.20/第4周} {
    \draw[BrandGray] (\x,2.08) -- (\x,2.32);
    \node[below] at (\x,2.02) {\t};
  }
  \node[draw=BrandBlue,fill=BrandBlueLight,rounded corners=2pt,text width=1.65cm,align=center,inner sep=5pt] at (1.1,3.25) {环境就绪\\doctor 通过};
  \node[draw=BrandPurple,fill=BrandPurpleLight,rounded corners=2pt,text width=1.65cm,align=center,inner sep=5pt] at (2.75,3.25) {专家 profile\\知识包};
  \node[draw=BrandOrange,fill=BrandOrangeLight,rounded corners=2pt,text width=1.65cm,align=center,inner sep=5pt] at (4.4,3.25) {业务试用\\问题集评测};
  \node[draw=BrandGreen,fill=BrandGreenLight,rounded corners=2pt,text width=1.65cm,align=center,inner sep=5pt] at (5.85,3.25) {试点报告\\扩展建议};
  \draw[-{Stealth[length=2mm]},BrandOrange,thick] (1.90,3.25) -- (2.00,3.25);
  \draw[-{Stealth[length=2mm]},BrandOrange,thick] (3.55,3.25) -- (3.70,3.25);
  \draw[-{Stealth[length=2mm]},BrandOrange,thick] (5.18,3.25) -- (5.28,3.25);
  \node[draw=BrandOrangeDark,fill=white,rounded corners=2pt,text width=5.8cm,align=center,inner sep=6pt,font=\bfseries] at (3.22,0.78) {最小决策：不是生产上线，而是批准一个可控、可审计、可量化的内网试点。};
\end{tikzpicture}
\end{column}
\end{columns}
\end{frame}

\end{document}
